\documentclass[11pt]{article}

\usepackage[margin=1in]{geometry}
\usepackage{amsmath,amssymb,bm}
\usepackage{booktabs}
\usepackage{array}
\usepackage{microtype}
\usepackage{hyperref}
\usepackage{xurl}
\usepackage[T1]{fontenc}
\usepackage{lmodern}

\hypersetup{
  colorlinks=true,
  linkcolor=blue,
  citecolor=blue,
  urlcolor=blue
}

\newcommand{\SFV}{SFV/dSB}
\newcommand{\PS}{Pati--Salam}
\newcommand{\MSbar}{\overline{\mathrm{MS}}}
\newcommand{\GitHubURL}{https://github.com/USERNAME/REPOSITORY}
\newcommand{\ZenodoDOI}{DOI-TO-BE-ASSIGNED}

\title{\textbf{Pati--Salam Seed Multiplicity and a Constrained Internal-Response Metric for Quark Flavor in SFV/dSB}\\[0.4em]
\large Phase B2 v1.7.0}
\author{Steven Hoffmann\\Independent Researcher}
\date{22 August 2026}

\begin{document}
\maketitle

\begin{abstract}
We document Phase B2 v1.7.0 of the SFV/dSB flavor program, a reproducible construction in which quark mass ratios and CKM magnitudes are generated by chiral fermion localization on a frozen corrected $O(4)$ scalar-wall background. The operator-origin chain is developed in several steps. A complete adjoint seed of $SU(4)_C\times SU(2)_L\times SU(2)_R$ supplies the multiplicity $15+3+3=21$. In a protected zero-momentum auxiliary matching, one singlet seed plus the 21-component adjoint produces the two-channel Gram matrix $\bigl(\begin{smallmatrix}1&1\\1&22\end{smallmatrix}\bigr)$, whose inverse yields the rational sector relations $22/21$, $23/21$, and $1/21$. The corrected bounce fixes an independent radial driver, $\Sigma_{\rm rad}=2(1-\Phi_c/\rho)^2(1+\Phi_c/\rho)=0.0670851132$. A minimal determinant-preserving internal-response action with isotropic blocks of dimensions 21 and 15 then uniquely gives the generator $T=\mathrm{diag}(\mathbf 1_{21}/21,-\mathbf 1_{15}/15)$ and the factors $Z_0=e^{\Sigma/21}$ and $Z_F=e^{-2\Sigma/15}$. With the architecture and frozen inputs fixed, the resulting seven-observable benchmark has a maximum relative discrepancy of $0.6111\%$ and an RMS discrepancy of $0.3580\%$, without adjusting a continuous flavor amplitude at the v1.7.0 stage. The preferred response fields are auxiliary and add neither a spacetime dimension nor a propagating scalar. We state a strict claim boundary: the compensator locking invariant is an additional structural principle beyond the original two-scalar SFV/dSB action, and the historical development inspected the flavor benchmark. The result is therefore a zero-continuous-fit realization and minimal effective completion, not an independently blind first-principles prediction of the unextended action.
\end{abstract}

\section{Introduction}

Localized chiral zero modes on solitonic backgrounds provide a standard mechanism for converting background-field structure into effective fermion hierarchies \cite{JackiwRebbi1976,RubakovShaposhnikov1983}. The SFV/dSB program applies this logic to a corrected two-field $O(4)$ tunneling background and asks a stronger question: once the wall background is frozen, can the coefficients entering the localized quark operator be traced to explicit symmetry, multiplicity, and response structures rather than treated as independent flavor-fit parameters?

Phase B2 addresses that operator-origin problem. The construction uses a Pati--Salam gauge organization \cite{PatiSalam1974}, a complete adjoint seed, a protected auxiliary response, and a bounce-derived radial normalization. The final v1.7.0 checkpoint introduces an explicit constrained internal-response-metric action that generates the two residual exponential compensators required by the frozen flavor architecture.

The principal numerical fact is simple: after the operator architecture and bounce inputs are fixed, seven quark-flavor observables are reproduced with a worst relative discrepancy of $0.6111\%$. The purpose of this paper is not merely to report that number, but to separate what is derived, what is protected, what is conditional, and what remains an added structural postulate.

A further methodological qualification is essential. Earlier stages of the program inspected the target flavor benchmark while searching for the successful architecture. Thus ``zero-continuous-fit'' below means that no continuous flavor coefficient is adjusted in the frozen v1.7.0 evaluation. It does \emph{not} mean that the complete architecture was selected in a blind holdout experiment.

\section{Frozen inputs and observable set}

The scalar background is treated as a frozen input to Phase B2. The benchmark uses the corrected regular $O(4)$ profile included in the repository at\\
\path{data/background_profile_O4_regular_robin_full.csv}.
No wall parameter is retuned in the v1.7.0 flavor evaluation.

The seven dimensionless observables are
\begin{equation}
\left\{
\frac{m_c}{m_t},\frac{m_u}{m_t},\frac{m_s}{m_b},\frac{m_d}{m_b},
|V_{us}|,|V_{cb}|,|V_{ub}|
\right\}.
\end{equation}
The frozen mass ratios are evaluated at $\mu=M_Z$ using the full Standard Model $\MSbar$ running masses tabulated by Huang and Zhou \cite{HuangZhou2021}. The repository records the numerical target vector and the CKM benchmark in \texttt{configs/targets\_MZ.json}.

The calculation is deterministic once the frozen architecture is specified. Relative errors reported below are benchmark discrepancies, not likelihood residuals or uncertainty-normalized pulls.

\section{Localized chiral operator}

The earlier Phase B2 operator audit begins from a Hermitian transverse fermion problem with brane-parallel and transverse kinetic weights and a wall-dependent mass kernel. After canonical normalization, the chiral first-order equations can be written schematically as
\begin{equation}
 f'_{L,R}
 =-\mathcal A_y f_{L,R}
 -\frac{1}{2}(\ln A)'f_{L,R}
 \mp A^{-1}M f_{L,R},
\end{equation}
where the sign of the mass term reverses between the two chiralities. A key correction in this stage was that a common isotropic position-dependent kinetic factor cancels after proper Hermitian normalization; an apparent scalar kinetic connection cannot by itself generate an independent diagonal flavor potential.

For the successful reduced wall operator, the localized profile is organized by an odd transition structure $O(y)$ and an even gradient-core structure $E(y)$,
\begin{equation}
 f'(y)=-\bigl[q\,O(y)+\kappa\,E(y)\bigr]f(y),
\end{equation}
with sector and family dependence residing in the protected coefficients. Spectral checks in the repository preserve nine desired chiral zero modes and exclude the nine opposite-chirality partners.

\section{Origin of the integer 21 from Pati--Salam seed multiplicity}

The Pati--Salam group is
\begin{equation}
G_{PS}=SU(4)_C\times SU(2)_L\times SU(2)_R.
\end{equation}
Its adjoint dimension is
\begin{equation}
\dim\mathrm{Adj}(G_{PS})=15+3+3=21.
\end{equation}
The relevant integer is therefore a full gauge-algebra multiplicity, not a fitted numerical approximation.

Introduce two real singlet mediator coordinates $S=(S_1,S_2)$, one heavy real singlet seed $X_0$, and a complete real Pati--Salam adjoint seed $X_{\mathcal A}$, $\mathcal A=1,\ldots,21$. A renormalizable local seed sector may be written
\begin{align}
\mathcal L_{\rm seed}={}&
\frac12(\partial X_0)^2+\frac12(DX_{\mathcal A})^2
-\frac12\left[M_X^2+2\mu(S_1+S_2)\right]X_0^2 \nonumber\\
&-\frac12\left[M_X^2+2\mu S_2\right]X_{\mathcal A}X_{\mathcal A}
+\mathcal L_{\rm stab}.
\end{align}
The corresponding incidence vectors are
\begin{equation}
v_0=(1,1),\qquad v_{\rm adj}=(0,1).
\end{equation}
For equal per-component zero-momentum response, the induced quadratic kernel is proportional to
\begin{align}
K_S
&=v_0v_0^T+\sum_{\mathcal A=1}^{21}v_{\rm adj}v_{\rm adj}^T\\
&=\begin{pmatrix}1&1\\1&1\end{pmatrix}
+21\begin{pmatrix}0&0\\0&1\end{pmatrix}
=\boxed{\begin{pmatrix}1&1\\1&22\end{pmatrix}}.
\end{align}
Its determinant is 21 and
\begin{equation}
K_S^{-1}=\frac1{21}
\begin{pmatrix}22&-1\\-1&1\end{pmatrix}.
\end{equation}

With the right-handed isospin routing $T_{3R}(Q_L)=0$, $T_{3R}(u_R)=+1/2$, and $T_{3R}(d_R)=-1/2$, contractions with the inverse kernel yield the frozen sector relations
\begin{equation}
\boxed{
\frac{h_Q}{h_{d0}}=\frac{22}{21},\qquad
\frac{h_{u0}}{h_{d0}}=\frac{23}{21},\qquad
\frac{h_{u1}}{h_{d1}}=\frac1{21}.}
\end{equation}
The family-odd direction is carried by the anomaly-free generator $F=\mathrm{diag}(-1,0,+1)$ in the frozen construction.

\section{Protected auxiliary matching}

A central consistency question is whether canonical normalization changes the seed-response ratio into a different loop kernel. Phase B2 therefore introduced a zero-momentum auxiliary/Legendre response description. Let $Y_0$ be conjugate to $X_0^2$ and $Y_{\rm Adj}$ to the normalized adjoint bilinear. The first-order response action has the schematic form
\begin{equation}
\Gamma_{\rm aux}
=\frac12\frac{M_0^2}{\mu_0}Y_0^2-Y_0(S_1+S_2)
+21\left[
\frac12\frac{M_{\rm Adj}^2}{\mu_{\rm Adj}}Y_{\rm Adj}^2-Y_{\rm Adj}S_2
\right].
\end{equation}
Elimination gives
\begin{equation}
K_S=\frac{\mu_0}{M_0^2}v_0v_0^T
+21\frac{\mu_{\rm Adj}}{M_{\rm Adj}^2}v_{\rm adj}v_{\rm adj}^T.
\end{equation}
At equal response $\mu_{\rm Adj}/M_{\rm Adj}^2=\mu_0/M_0^2$, the exact $21$ Gram structure follows.

The one-loop Pati--Salam gauge contribution obeys
\begin{equation}
(16\pi^2)\,\beta(M_a^2)|_g=-6C_2(\mathrm{adj}_a)g_a^2M_a^2,
\end{equation}
\begin{equation}
(16\pi^2)\,\beta(\mu_a)|_g=-6C_2(\mathrm{adj}_a)g_a^2\mu_a,
\end{equation}
so that
\begin{equation}
\frac{d}{d\ln Q}\ln\!\left(\frac{\mu_a}{M_a^2}\right)=0
\end{equation}
at one-loop gauge order. This protects the zero-momentum auxiliary response ratio even though $M_a^2$ and $\mu_a$ separately run.

This statement is specific to the protected auxiliary kernel. Other dynamical kernels need not be protected. In particular, the repository explicitly finds that a derivative-bubble scaling $\mu^2/M^2$ drives the effective multiplicity away from 21 and fails the sub-percent benchmark over the tested scale interval. That negative result is retained as a falsification check rather than discarded.

\section{Bounce-derived radial driver}

The corrected bounce center lies almost purely along the bulk radial soft mode. With true-vacuum radius $\rho$ and center value $\Phi_c$, define
\begin{equation}
x_c=\frac{\Phi_c}{\rho}=0.8659239806,
\qquad
u=1-x_c=0.1340760194.
\end{equation}
The normalized residual quartic energy is
\begin{equation}
\epsilon_c=(1-x_c^2)^2=0.0625878607.
\end{equation}
The finite-amplitude conversion from the harmonic radial normalization to the exact quartic displacement is
\begin{equation}
\lambda_{\rm radial}=\frac{1}{1-\nu/2}=1.0718550277.
\end{equation}
Their product simplifies to the closed form
\begin{equation}
\boxed{
\Sigma_{\rm rad}
=\lambda_{\rm radial}\epsilon_c
=2(1-x_c)^2(1+x_c)
=0.0670851132.}
\end{equation}
No flavor observable enters this radial number. The all-wall audit verifies the underlying residual-energy identity across 51 nearby wall solutions to numerical precision.

\section{Constrained internal-response metric}

The v1.7.0 completion introduces two isotropic response blocks,
\begin{equation}
C_{21}=e^a\mathbf 1_{21},\qquad
C_{15}=e^b\mathbf 1_{15}.
\end{equation}
The block dimensions have distinct origins: 21 is the Pati--Salam adjoint-seed multiplicity, while 15 is the fermion-source response block appearing in the frozen compensator architecture. The construction does not identify their sum with a physical spacetime dimension.

Impose the determinant constraint
\begin{equation}
21a+15b=0
\end{equation}
and define
\begin{equation}
\Sigma\equiv\ln\det C_{21}=21a.
\end{equation}
Block isotropy then fixes
\begin{equation}
\boxed{a=\frac{\Sigma}{21},\qquad b=-\frac{\Sigma}{15}},
\end{equation}
so the one-dimensional traceless generator is uniquely
\begin{equation}
\boxed{
T=\mathrm{diag}\!\left(
\frac1{21}\mathbf 1_{21},
-\frac1{15}\mathbf 1_{15}
\right)},
\end{equation}
with $\mathrm{Tr}\,T=0$ and $\mathrm{Tr}\,T^2=4/35$.

Write the logarithmic block volumes as
\begin{equation}
A=\ln\det C_{21},\qquad B=\ln\det C_{15}.
\end{equation}
The explicit auxiliary response potential is
\begin{equation}
\boxed{
V_{\rm resp}
=\frac{\Lambda_U^4}{2}(A+B)^2
+\frac{\Lambda_L^4}{2}
\left[\frac{A-B}{2}-\Sigma_{\rm rad}\right]^2
+V_{\rm shape}.}
\end{equation}
Here $V_{\rm shape}$ is positive and removes traceless distortions inside either isotropic block. For positive stiffnesses, the minimum is independent of the stiffness values:
\begin{equation}
A=\Sigma_{\rm rad},\qquad B=-\Sigma_{\rm rad}.
\end{equation}
Thus
\begin{equation}
\boxed{Z_0=e^{\Sigma_{\rm rad}/21}=1.0031996371},
\end{equation}
\begin{equation}
\boxed{Z_F=e^{-2\Sigma_{\rm rad}/15}=0.9910952029}.
\end{equation}
The factor of two in $Z_F$ follows from the two fermion legs in the source response.

In the preferred realization the response blocks are algebraic auxiliary variables without kinetic terms. They therefore add no propagating particle and no additional spacetime dimension. A kinetic term is optional and corresponds instead to a heavy propagating modulus.

\section{Frozen seven-observable result}

Table~\ref{tab:flavor} gives the frozen v1.7.0 targets, predictions, and relative discrepancies. The predictions are obtained from the repository target values and the deterministic v1.7.0 evaluation.

\begin{table}[ht]
\centering
\caption{Frozen Phase B2 v1.7.0 flavor benchmark.}
\label{tab:flavor}
\begin{tabular}{lccc}
\toprule
Observable & Target & Prediction & Relative error \\
\midrule
$m_c/m_t$ & $3.684774\times10^{-3}$ & $3.662257\times10^{-3}$ & $-0.6111\%$ \\
$m_u/m_t$ & $7.310115\times10^{-6}$ & $7.307888\times10^{-6}$ & $-0.0305\%$ \\
$m_s/m_b$ & $1.872490\times10^{-2}$ & $1.875711\times10^{-2}$ & $+0.1720\%$ \\
$m_d/m_b$ & $9.404720\times10^{-4}$ & $9.381294\times10^{-4}$ & $-0.2491\%$ \\
$|V_{us}|$ & $0.224300$ & $0.225095$ & $+0.3545\%$ \\
$|V_{cb}|$ & $0.042200$ & $0.0421402$ & $-0.1418\%$ \\
$|V_{ub}|$ & $0.003940$ & $0.00396104$ & $+0.5341\%$ \\
\bottomrule
\end{tabular}
\end{table}

The headline measures are
\begin{equation}
\boxed{\max_i |\delta_i|=0.6111\%},
\qquad
\boxed{\delta_{\rm RMS}=0.3580\%}.
\end{equation}
No continuous flavor amplitude is adjusted at the v1.7.0 stage.

This accuracy should not be confused with a seven-observable statistical goodness-of-fit. The target vector helped guide historical model development, and the theory uncertainties of the new effective sector have not been promoted to a full probabilistic error model. The appropriate interpretation is therefore structural consistency of a frozen realization.

\section{Robustness and propagating alternative}

The response locking is not numerically fine-tuned at the sub-percent boundary. Under a common rescaling
\begin{equation}
\Sigma\rightarrow r\,\Sigma_{\rm rad},
\end{equation}
all seven tested discrepancies remain below $1\%$ for approximately
\begin{equation}
0.706\lesssim r\lesssim1.361.
\end{equation}
Independent seed-only and fermion-only response scans allow still broader intervals around the central construction.

If the relative response mode is made propagating, a simple finite-momentum response model is
\begin{equation}
r(p)=\frac{m_\Sigma^2}{m_\Sigma^2+p^2}.
\end{equation}
The numerical tolerance scan gives
\begin{equation}
\boxed{m_\Sigma/p\gtrsim1.55}
\end{equation}
for the seven flavor discrepancies to remain below $1\%$. The auxiliary realization avoids this condition entirely.

At one-loop gauge order, the determinant constraint is algebraic and exact, the protected seed ratio $\mu/M^2$ is invariant, and multiplicative fermion-source running cancels out of the normalized radial exponent. A generic nonlinear source running can eventually renormalize the effective charge, but the v1.7.0 scan shows that corrections must become strong before reaching the direct flavor tolerance boundary.

\section{Claim boundary and falsifiability}

The scientific status can be stated compactly.

\subsection*{Derived or exact within the augmented construction}
\begin{itemize}
\item The Pati--Salam adjoint multiplicity is exactly $15+3+3=21$ once one complete adjoint seed is specified.
\item The protected auxiliary response with equal per-component $\mu/M^2$ gives the exact Gram matrix $\bigl(\begin{smallmatrix}1&1\\1&22\end{smallmatrix}\bigr)$ and its rational inverse relations.
\item The one-loop gauge contribution protects $\mu/M^2$ in the explicit seed-sector audit.
\item The corrected bounce fixes $\epsilon_c$, $\lambda_{\rm radial}$, and $\Sigma_{\rm rad}$ without using flavor data.
\item Given block dimensions 21 and 15, block isotropy, determinant one, and the convention $\Sigma=\ln\det C_{21}$, the generator $T$ is unique.
\item The explicit response potential has the exact minimum $A=\Sigma_{\rm rad}$ and $B=-\Sigma_{\rm rad}$.
\item With the frozen architecture, the downstream seven-observable realization has $0.6111\%$ maximum relative discrepancy and passes the repository regression suite.
\end{itemize}

\subsection*{Conditional or additional structure}
\begin{itemize}
\item Pati--Salam symmetry does not by itself require the UV theory to contain exactly one complete adjoint seed plus the singlet seed used here.
\item The determinant-preserving response compensator and, especially, its locking to the bounce radial driver are additional structural principles beyond the original two-scalar SFV/dSB Lagrangian.
\item A fully dynamical composite realization would have to reproduce the protected zero-momentum response and control finite-momentum derivative nonuniversality.
\item The historical model-development path was not blind to the seven flavor targets; future external holdouts are needed to convert structural success into genuinely prospective predictive evidence.
\end{itemize}

The construction is therefore falsifiable in several directions. A UV completion that necessarily generates an unprotected kernel, a response metric incompatible with the determinant constraint, a propagating modulus too light to satisfy the finite-momentum bound, or external flavor observables inconsistent with the frozen operator architecture would each weaken or exclude this particular completion.

\section{Reproducibility and archive structure}

The publication repository contains the frozen input tables, all major intermediate checkpoint reports, the v1.7.0 source code, numerical result files, and regression tests. The core reproduction commands are
\begin{verbatim}
python -m pip install -r requirements.txt
python src/constrained_internal_response_metric.py
pytest -q
\end{verbatim}
In the publication audit, the test suite returns 74 passing tests. The main machine-readable summary is
\begin{verbatim}
results/constrained_internal_response_metric/
  constrained_internal_response_metric_summary.json
\end{verbatim}
The earlier Phase B v0.1.0 microphysical-closure repository is preserved under \texttt{historical/} for provenance. It is not required to reproduce v1.7.0 and is explicitly labeled as an earlier partial/post-hoc closure stage.

The intended public repository location is \url{\GitHubURL}. The archival DOI will be inserted here after Zenodo publication: \texttt{\ZenodoDOI}.

\section{Conclusion}

Phase B2 v1.7.0 replaces a large part of the earlier empirical flavor-control structure with a concrete operator-origin chain. The Pati--Salam adjoint supplies the integer 21; a protected auxiliary response turns that multiplicity into the exact two-channel Gram matrix; the corrected bounce supplies a closed radial driver; and a minimal determinant-preserving response action fixes the residual exponents with no continuous flavor coefficient adjusted in the final evaluation. The resulting seven-observable benchmark remains within $0.6111\%$ in relative discrepancy.

The result is strongest when described narrowly. It is not yet a theorem that the original two-scalar SFV/dSB action uniquely implies the full quark-flavor sector. Rather, it is an explicit, reproducible, internally consistent minimal completion whose remaining new ingredient is sharply isolated: the response-compensator locking principle. That isolation is useful scientifically because future work can now target one structural question rather than a collection of unconstrained flavor coefficients.

\begin{thebibliography}{9}

\bibitem{PatiSalam1974}
J.~C. Pati and A.~Salam,
``Lepton number as the fourth color,''
\emph{Phys. Rev. D} \textbf{10}, 275--289 (1974),
\href{https://doi.org/10.1103/PhysRevD.10.275}{doi:10.1103/PhysRevD.10.275}.

\bibitem{JackiwRebbi1976}
R.~Jackiw and C.~Rebbi,
``Solitons with fermion number 1/2,''
\emph{Phys. Rev. D} \textbf{13}, 3398--3409 (1976),
\href{https://doi.org/10.1103/PhysRevD.13.3398}{doi:10.1103/PhysRevD.13.3398}.

\bibitem{RubakovShaposhnikov1983}
V.~A. Rubakov and M.~E. Shaposhnikov,
``Do we live inside a domain wall?,''
\emph{Phys. Lett. B} \textbf{125}, 136--138 (1983),
\href{https://doi.org/10.1016/0370-2693(83)91253-4}{doi:10.1016/0370-2693(83)91253-4}.

\bibitem{HuangZhou2021}
G.-y. Huang and S.~Zhou,
``Precise values of running quark and lepton masses in the Standard Model,''
\emph{Phys. Rev. D} \textbf{103}, 016010 (2021),
\href{https://arxiv.org/abs/2009.04851}{arXiv:2009.04851}.

\end{thebibliography}

\end{document}
